\section{Coin Change II (Count the Ways)}
\label{sec:dp-coin-change-2}

\problemstatement{Given an array $\textit{coins}$ of denominations (each
in unlimited supply) and a target $\textit{amount}$, count the
\textbf{number of distinct combinations} of coins that sum to exactly
$\textit{amount}$.}

\begin{patternbox}
This is Section~\ref{sec:dp-coin-change}'s recurrence with minimum
replaced by \textbf{addition} -- the same exists-vs-count distinction
drawn throughout this chapter and the Recursion chapter before it, now
applied to the unbounded-reuse setting.
\end{patternbox}

\subsection*{Deriving the recurrence}

\begin{gather*}
  f(\textit{index}, \textit{amount}) = f(\textit{index}{-}1, \textit{amount})
  \ +\ \textit{take}, \\
  \textit{take} = \textit{coins}[\textit{index}] \le \textit{amount}\ ?\
  f(\textit{index}, \textit{amount}-\textit{coins}[\textit{index}]) : 0.
\end{gather*}
Base case: $f(0, \textit{amount}) = 1$ if $\textit{amount}$ is exactly
divisible by $\textit{coins}[0]$ (exactly one way: use that many copies
of the first coin), else $0$.

\subsection*{Approach 1: Recursion}

\begin{lstlisting}
#include <bits/stdc++.h>
using namespace std;

int changeRecursive(int index, int amount, vector<int>& coins) {
    if (index == 0) {
        return (amount % coins[0] == 0) ? 1 : 0;
    }

    int notTake = changeRecursive(index - 1, amount, coins);
    int take = 0;
    if (coins[index] <= amount) {
        take = changeRecursive(index, amount - coins[index], coins);
    }
    return notTake + take;
}

int main() {
    vector<int> coins = {1, 2, 5};
    int n = coins.size();
    cout << changeRecursive(n - 1, 5, coins) << endl; // 4
    return 0;
}
\end{lstlisting}

\begin{complexitybox}[Approach 1 Complexity]
\textbf{Time.} Exponential, same overlapping-subproblem blowup as
Section~\ref{sec:dp-coin-change}.
\textbf{Space.} $O(\textit{amount})$ recursion stack.
\end{complexitybox}

\subsection*{Approach 2: Memoization}

\begin{lstlisting}
#include <bits/stdc++.h>
using namespace std;

int changeMemo(int index, int amount, vector<int>& coins,
                vector<vector<int>>& memo) {
    if (index == 0) {
        return (amount % coins[0] == 0) ? 1 : 0;
    }
    if (memo[index][amount] != -1) return memo[index][amount];

    int notTake = changeMemo(index - 1, amount, coins, memo);
    int take = 0;
    if (coins[index] <= amount) {
        take = changeMemo(index, amount - coins[index], coins, memo);
    }
    return memo[index][amount] = notTake + take;
}

int main() {
    vector<int> coins = {1, 2, 5};
    int n = coins.size(), amount = 5;
    vector<vector<int>> memo(n, vector<int>(amount + 1, -1));
    cout << changeMemo(n - 1, amount, coins, memo) << endl; // 4
    return 0;
}
\end{lstlisting}

\begin{complexitybox}[Approach 2 Complexity]
\textbf{Time.} $O(n \cdot \textit{amount})$.
\textbf{Space.} $O(n \cdot \textit{amount})$ memo $+$
$O(\textit{amount})$ recursion stack.
\end{complexitybox}

\subsection*{Approach 3: Tabulation}

\begin{lstlisting}
#include <bits/stdc++.h>
using namespace std;

int changeTabulation(vector<int>& coins, int amount) {
    int n = coins.size();
    vector<vector<long long>> dp(n, vector<long long>(amount + 1, 0));

    for (int a = 0; a <= amount; a++) {
        dp[0][a] = (a % coins[0] == 0) ? 1 : 0;
    }

    for (int index = 1; index < n; index++) {
        for (int a = 0; a <= amount; a++) {
            long long notTake = dp[index - 1][a];
            long long take = (coins[index] <= a) ? dp[index][a - coins[index]] : 0;
            dp[index][a] = notTake + take;
        }
    }
    return dp[n - 1][amount];
}

int main() {
    vector<int> coins = {1, 2, 5};
    cout << changeTabulation(coins, 5) << endl; // 4
    return 0;
}
\end{lstlisting}

\subsection*{Dry run of the tabulation table}

\dryrun{Take $\textit{coins}=[1,2,5]$, $\textit{amount}=5$.}

The four ways to make $5$: $\{1,1,1,1,1\}$, $\{1,1,1,2\}$,
$\{1,2,2\}$, $\{5\}$ -- matching $\textit{dp}[2][5]=4$. Row $0$ (coin
$1$ only): $\textit{dp}[0][a]=1$ for every $a$ (always exactly one way
using only $1$s). Row $1$ (coins $1,2$): combinations using some
$2$s and the rest $1$s -- $\textit{dp}[1][5]=3$ (zero, one, or two
$2$s, with the remainder in $1$s). Row $2$ adds the option of using
coin $5$ once, contributing the fourth way.

\begin{complexitybox}[Approach 3 Complexity]
\textbf{Time.} $O(n \cdot \textit{amount})$.
\textbf{Space.} $O(n \cdot \textit{amount})$ table, no recursion stack.
\end{complexitybox}

\subsection*{Approach 4: Space Optimization}

\begin{lstlisting}
#include <bits/stdc++.h>
using namespace std;

int changeSpaceOptimized(vector<int>& coins, int amount) {
    int n = coins.size();
    vector<long long> prev(amount + 1, 0);

    for (int a = 0; a <= amount; a++) {
        prev[a] = (a % coins[0] == 0) ? 1 : 0;
    }

    for (int index = 1; index < n; index++) {
        vector<long long> current(amount + 1, 0);
        for (int a = 0; a <= amount; a++) {
            long long notTake = prev[a];
            long long take = (coins[index] <= a) ? current[a - coins[index]] : 0;
            current[a] = notTake + take;
        }
        prev = current;
    }
    return prev[amount];
}

int main() {
    vector<int> coins = {1, 2, 5};
    cout << changeSpaceOptimized(coins, 5) << endl; // 4
    return 0;
}
\end{lstlisting}

\begin{complexitybox}[Approach 4 Complexity]
\textbf{Time.} $O(n \cdot \textit{amount})$.
\textbf{Space.} $O(\textit{amount})$ for the rolling row, reading
``take'' from \textit{current} exactly as in
Section~\ref{sec:dp-coin-change}.
\end{complexitybox}

\begin{takeawaybox}
Coin Change and Coin Change II are, structurally, the exact same
recurrence with minimum versus addition as the only difference -- the
unbounded-knapsack equivalent of Subset Sum versus Count Subsets earlier
in this chapter. Once you have one of a min/count/exists triplet for a
given recurrence shape, the other two cost almost nothing extra to
derive.
\end{takeawaybox}
